% !TeX root = ../navier-stokes-manuscript.tex
\subsection{Why the finite pieces are rectangles}\label{sub:BodyFiniteAdjacentRectangles}

Section~\ref{sec:CriticalHeightAndDerivativeTower} ended with a finite block of the mixed grid.
Let \((u,p)\) be the maximal classical solution generated by the fixed pair \((u_0,\nu)\), and write its lifespan as \([0,T_*)\).
On the finite-lifespan branch used below, set
\[
  I_*:=(0,T_*).
\]
To turn its spatial coordinates into time traces, pair each response \(V_n\) with its time derivative \(V_{n+1}\).
Let
\(A=(1-\Delta)^{1/2}\).  For a smooth Hilbert-valued function \(v\) and
\(0\leq a<b\leq T\),
\begin{equation}\label{eq:BodyAdjacentTraceIdentity}
  \norm{v(b)}_{H^M}^2-
  \norm{v(a)}_{H^M}^2
  =2\operatorname{Re}
  \int_a^b
  \pair{A^{M-1}v_t(t)}{A^{M+1}v(t)}_{L^2}\,dt.
\end{equation}
Cauchy--Schwarz now places the change of the middle norm between its two adjacent
space-time rows:
\[
  \left|
  \norm{v(b)}_{H^M}^2-
  \norm{v(a)}_{H^M}^2
  \right|
  \leq
  2\norm{v_t}_{L^2(a,b;H^{M-1})}
   \norm{v}_{L^2(a,b;H^{M+1})}.
\]
Consequently,
\begin{equation}\label{eq:BodyAdjacentTracePair}
  v\in L^2(0,T;H^{M+1}),
  \qquad
  v_t\in L^2(0,T;H^{M-1})
\end{equation}
gives
\begin{equation}\label{eq:BodyAdjacentTraceConclusion}
  v\in C([0,T];H^M).
\end{equation}
Time mollification and Fourier truncation justify
\eqref{eq:BodyAdjacentTraceIdentity} for the pair in
\eqref{eq:BodyAdjacentTracePair}; weak \(H^M\) continuity together with continuity
of the norm gives the strong representative.

For \(M=2\), the pair is already concrete:
\begin{equation}\label{eq:FirstConcreteAdjacentPair}
  u\in L^2(0,T;H^3),
  \qquad
  u_t\in L^2(0,T;H^1)
  \quad\Longrightarrow\quad
  u\in C([0,T];H^2).
\end{equation}
The upper row reaches one derivative above the desired middle space.  The next
temporal row reaches one derivative below it.  Their pairing carries \(u\) through the
middle.

Apply this to \(v=V_n\), with \(v_t=V_{n+1}\).  If we want the first \(N+1\)
temporal responses to arrive in \(H^M\), the squared velocity norm used to measure that finite block is
\begin{equation}\label{eq:BodyRectangleNorm}
\begin{aligned}
  \mathfrak R_{N,M}(u;T)
  :={}&
  \sum_{n=0}^{N}
  \norm{V_n}_{L^2(0,T;H^{M+1})}^2
  \\
  &+
  \sum_{n=0}^{N}
  \norm{V_{n+1}}_{L^2(0,T;H^{M-1})}^2.
\end{aligned}
\end{equation}
For \(M=0\), the lower edge is read in \(H^{-1}\).  Temporal depth \(N\) and
spatial height \(M\) remain independent.

\Needspace{12\baselineskip}
For example, \((N,M)=(1,2)\) contains two overlapping trace pairs:
\begin{center}
\captionsetup{hypcap=false}
\captionof{table}{The \((1,2)\) finite rectangle.  The row \(V_1\) is shared: it is the time derivative used to trace \(V_0\), and it is also the next response whose own trace is paired with \(V_2\).}
\label{tab:FirstAdjacentRectangle}
\renewcommand{\arraystretch}{1.35}
\begin{tabular}{@{}c c c@{}}
\toprule
temporal row & one step above & one step below \\
\midrule
\(n=0\) & \(V_0=u\in L^2H^3\) & \(V_1=u_t\in L^2H^1\) \\
\(n=1\) & \(V_1=u_t\in L^2H^3\) & \(V_2=u_{tt}\in L^2H^1\) \\
\bottomrule
\end{tabular}
\end{center}

Write \(\cR_{N,M}[u_0;T]\) for this finite pressure-complete block of the mixed jet:
the velocity coordinates measured by \eqref{eq:BodyRectangleNorm}, together with the
corresponding \(Q_n\)-coordinates reconstructed by
\eqref{eq:BodyMixedPressureRecurrence}.  Every entry is a coordinate of the
datum-generated velocity and pressure fields.

\divider{From the mixed-jet matrix to the Sobolev rectangles}

We can now say exactly what the matrix form records.
The mixed recurrence in \eqref{eq:BodyMixedJetRecurrence} and the Sobolev grid in Section~\ref{sub:ThePressureCompleteMixedDerivativeTower} are two readings of the same derivative family, but they do different jobs.
The first records how one differentiated transport row is assembled from the mixed jet.
Let
\[
  \Gamma:=\Nzero\times\Nzero^3,
  \qquad
  \kappa=(n,\alpha),
  \qquad
  \lambda=(a,\beta),
  \qquad
  \binom{\kappa}{\lambda}
  :=\binom{n}{a}\binom{\alpha}{\beta},
\]
and write
\[
  \lambda\preceq\kappa
  \quad\Longleftrightarrow\quad
  a\leq n\ \text{ and }\ \beta\leq\alpha,
  \qquad
  \kappa-\lambda=(n-a,\alpha-\beta).
\]
Set \(J_\kappa:=J_{n,\alpha}\) and
\(\mathbf J:=(J_\lambda)_{\lambda\in\Gamma}\).
The transport blocks are
\[
  \mathbb L(J)_{\kappa,\lambda}
  :=
  \binom{\kappa}{\lambda}
  \Leray\bigl[(J_{\kappa-\lambda}\cdot\nabla)\,\mathord\cdot\,\bigr],
  \qquad \lambda\preceq\kappa,
\]
with zero blocks when \(\lambda\npreceq\kappa\).
Their action is not left implicit.  The \(\kappa\)-th block row gives
\begin{align*}
  [\mathbb L(J)\mathbf J]_\kappa
  &=
  \sum_{\lambda\preceq\kappa}
  \binom{\kappa}{\lambda}
  \Leray\bigl[(J_{\kappa-\lambda}\cdot\nabla)J_\lambda\bigr]
  \\
  &=
  \sum_{\lambda\preceq\kappa}
  \binom{\kappa}{\lambda}
  \Leray\bigl[(J_\lambda\cdot\nabla)J_{\kappa-\lambda}\bigr].
\end{align*}
The second line is the change of index
\(\lambda\leftrightarrow\kappa-\lambda\), using
\(\binom{\kappa}{\kappa-\lambda}=\binom{\kappa}{\lambda}\).
It is exactly the Leray projection of the nonlinear sum in
\eqref{eq:BodyMixedJetRecurrence}, so the projected velocity row is
\[
  J_{n+1,\alpha}
  +[\mathbb L(J)\mathbf J]_\kappa
  =\nu\Delta J_{n,\alpha}.
\]
The projection removes \(\nabla\Pi_{n,\alpha}\) from this velocity row because
\(\Leray\nabla\Pi_{n,\alpha}=0\); the pressure coordinate remains reconstructed by
\eqref{eq:BodyMixedPressureRecurrence} and retained in the pressure-complete rectangle.

Order \(\Gamma\) by the total mixed depth \(n+|\alpha|\), with any fixed rule for ties.
The condition \(\lambda\preceq\kappa\) then places the supported blocks on or below the diagonal.
This triangularity belongs to the transport bookkeeping: the endpoint terms
\(\lambda=0\) and \(\lambda=\kappa\) still contain \(J_\kappa\), while the full equation also contains the forward temporal shift and the Laplacian.
Each block is a first-order transport operator followed by the order-zero Leray projection; it depends on the solution and is generally nonlocal.

The Pascal weights count how the derivatives divide between the two velocity factors.
They disappear if each coordinate is measured with its own factorial.
With \(\alpha!:=\alpha_1!\alpha_2!\alpha_3!\), set
\[
  \widehat J_{n,\alpha}
  :=\frac{J_{n,\alpha}}{n!\,\alpha!},
  \qquad
  \widehat\Pi_{n,\alpha}
  :=\frac{\Pi_{n,\alpha}}{n!\,\alpha!}.
\]
The cancellation and the surviving temporal factor can be seen in one summand:
\begin{align*}
  \frac{J_{n+1,\alpha}}{n!\,\alpha!}
  &=(n+1)\widehat J_{n+1,\alpha},
  \\
  \frac{1}{n!\,\alpha!}
  \binom{n}{a}\binom{\alpha}{\beta}
  (J_{a,\beta}\cdot\nabla)J_{n-a,\alpha-\beta}
  &=(\widehat J_{a,\beta}\cdot\nabla)
    \widehat J_{n-a,\alpha-\beta}.
\end{align*}
Indeed,
\(\binom{n}{a}a!(n-a)!=n!\), and the same componentwise identity holds for
\(\binom{\alpha}{\beta}\beta!(\alpha-\beta)!=\alpha!\).
Dividing \eqref{eq:BodyMixedJetRecurrence} by \(n!\,\alpha!\) now gives
\[
  (n+1)\widehat J_{n+1,\alpha}
  +\sum_{a=0}^{n}\sum_{\beta\leq\alpha}
    (\widehat J_{a,\beta}\cdot\nabla)
    \widehat J_{n-a,\alpha-\beta}
  +\nabla\widehat\Pi_{n,\alpha}
  =\nu\Delta\widehat J_{n,\alpha}.
\]
Let \(D\) be the diagonal factorial matrix
\(D_{\kappa,\kappa}:=\kappa!:=n!\,\alpha!\).
Row by row---equivalently, on every finite truncation---the same cancellation reads
\[
  D^{-1}\mathbb L(D\widehat J)D
  =\widehat{\mathbb L}(\widehat J),
  \qquad
  \widehat{\mathbb L}(\widehat J)_{\kappa,\lambda}
  :=\Leray\bigl[(\widehat J_{\kappa-\lambda}\cdot\nabla)\,\mathord\cdot\,\bigr],
  \qquad \lambda\preceq\kappa.
\]
Each normalized block now depends on the difference \(\kappa-\lambda\), and every block row contains finitely many terms.
This is the lower-triangular Toeplitz, or convolution, form of the transport array.
The factor \(n+1\) remains on the temporal shift, and pressure remains attached through \(\widehat\Pi_{n,\alpha}\); the normalization has reorganized the original equation, not replaced it with a linear one.

The Sobolev grid gives the second matrix reading by measuring this same jet rather than listing each spatial multi-index separately.
For \(m\in\mathbb Z\), write
\[
  G_{n,m}(T)
  :=\norm{V_n}_{L^2(0,T;H^m)}
  =\left(\int_0^T\norm{A^mV_n(t)}_2^2\,dt\right)^{1/2}.
\]
Since \(J_{n,\alpha}=\partial_x^\alpha V_n\), the integer Sobolev norm satisfies
\[
  G_{n,m}(T)^2
  \asymp_m
  \sum_{|\alpha|\leq m}
  \norm{J_{n,\alpha}}_{L^2((0,T)\times\R^3)}^2,
  \qquad m\in\Nzero.
\]
Thus \(\alpha\) labels one spatial derivative in the mixed-jet array, whereas \(m\) labels the normed envelope that collects the spatial derivatives through order \(m\).
With this notation, \eqref{eq:BodyRectangleNorm} is exactly
\[
  \mathfrak R_{N,M}(u;T)
  =\sum_{n=0}^{N}G_{n,M+1}(T)^2
   +\sum_{n=0}^{N}G_{n+1,M-1}(T)^2.
\]
The rectangle is a finite norm window of the jet; it is not a principal submatrix of \(\mathbb L(J)\), nor does it add every entry beneath one corner.
For each \(0\leq n\leq N\), it selects the staggered trace stencil
\[
  (n,M+1),
  \qquad
  (n+1,M-1)
  \quad\longrightarrow\quad
  (n,M).
\]
The two input coordinates are the adjacent rows in \eqref{eq:BodyRectangleNorm}, and the output coordinate is the continuous middle trace produced by \eqref{eq:BodyAdjacentTraceIdentity}.
One step down the temporal direction spans two spatial indices and lands between them.
Increasing \(N\) stacks more of these stencils; changing \(M\) moves the same trace mechanism across the spatial columns.
The critical-height calculation in \eqref{eq:BodyCompletedCriticalRow} follows one spatial ladder inside the original response \(V_0=u\).
Equation~\eqref{eq:BodyEveryTemporalRowCompletedHeight} repeats that calculation inside each \(V_r\), opening the corresponding ladder in every temporal row, while the rectangles keep temporal depth and spatial height independent.

The operator array records how the mixed derivatives interact; the Sobolev grid records the sizes of those same derivatives in finite norm windows.
Both constructions come from the original pressure-complete history, but they are not the same matrix under a relabelling.
When a larger window contains a smaller one, a retained coordinate is the same distributional derivative of the same velocity or normalized pressure, viewed through the lower Sobolev height.
That identity, rather than the Toeplitz pattern by itself, is what makes restriction of the finite windows exact.
The rectangle theorem is a statement about these normed derivative coordinates, not about the operator norm, spectrum, or invertibility of \(\mathbb L(J)\).
The matrix form turns the Leibniz bookkeeping into a lower-triangular convolution with finitely many terms in each mixed-index row and identifies exactly which finite coordinates enter each trace.
It proves no estimate for any coordinate and improves no Sobolev norm.
The whole-terminal bound is the separate statement realized in \Cref{thm:WholeTerminalCompatibleRectangleEstimate}; no sum over all derivative orders and no constant shared by every rectangle is needed.

For any strict cutoff \(T<T_*\), classical smoothness makes every fixed block finite:
\begin{equation}\label{eq:StrictSubintervalRectangleBound}
  \mathfrak R_{N,M}(u;T)<\infty
  \qquad\text{for finite }N,M.
\end{equation}
Classicality supplies only the statement on the left:
\[
  \forall\,T<T_*:\quad \mathfrak R_{N,M}(u;T)<\infty
  \quad\not\Longrightarrow\quad
  \sup_{0<T<T_*}\mathfrak R_{N,M}(u;T)<\infty.
\]
The next subsection states the right-hand requirement as \Cref{thm:BodyWholeTerminalRectangleTheorem}: after \(N\) and \(M\) have been fixed, one bound must survive every cutoff.
